\documentclass[10pt,a4paper]{amsart}
\usepackage[margin=19mm]{geometry}
\usepackage{amsmath,amssymb,mathtools}
\usepackage[scr=rsfs,cal=euler]{mathalfa}
\usepackage{xfrac}
\usepackage[unicode,hidelinks,bookmarks=false]{hyperref}
\newcommand{\ie}{i.e.\ }
\newcommand{\cf}{cf.\ }
\newcommand{\resp}{respectively\ }
\newcommand{\Subsection}{\subsection}
\providecommand{\nobreakdash}{\nobreak\hbox{-}}
\setlength{\emergencystretch}{2em}
\multlinegap0pt
\usepackage{xcolor}
\usepackage{mathtools}
\usepackage{amssymb}
\usepackage{tikz}
\newtheorem{theorem}{Theorem}
\newtheorem{remark}{Remark}
\newtheorem{lemma}{Lemma}
\newcommand{\N}{\mathbb{N}}
\newcommand{\R}{\mathbb{R}}
\DeclareMathOperator{\Sep}{Sep}
\newcommand{\eqdef}{\vcentcolon=}
\title{On finiteness properties of the separating semigroup of a real curve}
\author{Matthew Magin}
\date{Source: arXiv:2511.18545v5 (CC BY 4.0)}
\begin{document}
\maketitle

\noindent\emph{Source and licence.} On finiteness properties of the separating semigroup of a real curve. Version arXiv:2511.18545v5 (CC BY 4.0), distributed by arXiv under the Creative Commons Attribution 4.0 International licence, http://creativecommons.org/licenses/by/4.0/, as recorded in the arXiv licence metadata for this exact version and shown on the abstract page https://arxiv.org/abs/2511.18545v5. Full licence notice: \texttt{licenses/arxiv-2511.18545-CC-BY-4.0.txt}.\par\smallskip

\noindent\textit{Editorial scope.} Counted unit: the article's theorem main\_thm (source0.tex lines 614--616). The article's complete original proof of it is delivered with it (source0.tex lines 617--645, one \texttt{proof} environment of 29 lines). No mathematics is written, repaired or re-derived: the statement and the proof are the article's own lines, in the article's own notation, and no retained line is substituted or re-flowed.  Retained uncounted as the source-local context the proof needs: lines 256--258; lines 395--397; lines 494--505.  Nothing was omitted from any retained span, so the package carries no omission record.  No retained line contains a citation, so the wrapper carries no bibliography and no bibliography entry.  No retained line contains a figure environment, an image-inclusion command or drawing code, so no image is delivered and the package declares no asset dependency.  Editorial formatting only, in addition: an AMS \texttt{amsart} preamble that restates the article's own declarations and macros used by the retained lines, namely the environments \texttt{theorem}, \texttt{remark}, \texttt{lemma} on separate counters, together with the standard packages those lines need.  The retained environments print in this document as Theorem 1, Remark 1, Lemma 1 in order of appearance, whereas the article prints the same items with the numbering of its own full document.  The complete native source of the article as distributed in the arXiv e-print for this exact version is bundled unchanged as \texttt{sources/189663/source0.tex}.
  \begin{theorem}\label{deleting_points}
      Let $P = p_1 + \ldots + p_n$ be a separating divisor on a real curve $X$ and suppose $n \geqslant g+2$. Then there exists a proper subset  $\Gamma \subsetneq \{ 1, \ldots, n\}$ such that the divisor $Q \eqdef \sum_{i \in \Gamma} p_i$ is separating and satisfies $\deg{Q} \geqslant \lceil \frac{n}{2} \rceil$. 
  \end{theorem}

  \begin{remark}\label{degree_g+2}
    Note that by the Riemann--Roch theorem, the hypothesis $h^0(P)\geqslant3$ of Theorem~\ref{deleting_points} holds whenever $\deg P\geqslant g+2$.
  \end{remark}

  \begin{lemma}\label{main_lemma}
  The separating semigroup of any real curve $X$ of genus $g$ with $b_0(\R X) = r$ admits the decomposition
  \[
    \Sep(X) = \Sep_s(X) \,\cup\, \bigcup_{i=1}^{m} \bigl( d(P_i) + \N_0^{\,r} \bigr),
  \]
  where $\Sep_s(X)$ is the finite set of degree partitions of special separating divisors, and $d(P_i)$ are the minimal (with respect to the pointwise partial order $\preceq$) degree partitions  corresponding to non-special separating divisors $P_i$.

  Moreover, we have
  \[
    |\Sep_s(X)| \leqslant \binom{2g - 3}{r} + \binom{g - 2}{r - 1}.
  \]
  \end{lemma}

  \begin{theorem}\label{main_thm}
    For a non-negative integer number $g$ the set of all separating semigroups of genus $g$ curves is finite. 
  \end{theorem}

  \begin{proof} 
    We may assume that $g \geqslant 2$, as the result for $g \leqslant 1$ follows from the known classification (see the introduction). By Harnack's inequality, any real curve $X$ of genus $g$ satisfies $b_0(\R X) \leqslant g + 1$.  Hence, it suffices to prove the claim for fixed $g$ and $r = b_0(\R X)$.

    Let $X$ be a real curve of genus $g$ with $b_0(\R X) = r$.  By Lemma~\ref{main_lemma}, its separating semigroup admits the decomposition
    \[
      \Sep(X) = \Sep_s(X) \,\cup\, \bigcup_{i=1}^{m} \bigl( d(P_i) + \N_0^{\,r} \bigr).
    \]

    Since for every $d \in \Sep_s(X)$ we have $|d| \leqslant 2g - 2$, there are only finitely many possibilities for $\Sep_s(X)$. Thus, in order to prove the theorem, it suffices to obtain a uniform bound on the degrees of the minimal non-special divisors $P_1, \ldots, P_m$.

    \begin{lemma}\label{4g-3_lemma}
    For $g \geqslant 2$, one has $\deg P_i \leqslant 4g - 3$.
    \end{lemma}

    \begin{proof}[Proof of Lemma~\ref{4g-3_lemma}]
    Suppose, for a contradiction, that $\deg P_i \geqslant 4g - 2$. 
    Since $g \geqslant 2$, we have $4g - 2 \geqslant g + 2$. 
    Then, by Theorem~\ref{deleting_points} (see Remark~\ref{degree_g+2}), there exists a separating divisor $Q_i$ such that $Q_i \leqslant P_i$, $Q_i \ne P_i$, and
    \[
      \deg Q_i \geqslant \left\lceil \frac{\deg P_i}{2} \right\rceil \geqslant 2g - 1.
    \]
    Hence, $Q_i$ is non-special (as $\deg Q_i \geqslant 2g - 1 > 2g - 2$) and separating, contradicting the minimality of $P_i$.
    \end{proof}

    It follows from Lemma~\ref{4g-3_lemma} that the number $m$ of minimal non-special divisors satisfies 
    \[m \leqslant (4g - 3)^r.\]
      Consequently, for fixed $g$ and $r$, there are only finitely many possible degree partitions $d(P_i)$, which implies the claim of Theorem~\ref{main_thm}.

    \end{proof}

\end{document}
